\documentclass[10pt,a4paper]{amsart}
\usepackage[margin=19mm]{geometry}
\usepackage{amsmath,amssymb,mathtools}
\usepackage[scr=rsfs,cal=euler]{mathalfa}
\usepackage{xfrac}
\usepackage[unicode,hidelinks,bookmarks=false]{hyperref}
\newcommand{\ie}{i.e.\ }
\newcommand{\cf}{cf.\ }
\newcommand{\resp}{respectively\ }
\newcommand{\Subsection}{\subsection}
\providecommand{\nobreakdash}{\nobreak\hbox{-}}
\setlength{\emergencystretch}{2em}
\multlinegap0pt
\usepackage{amssymb}
\usepackage[normalem]{ulem}
\usepackage{tikz}
\newtheorem{lemma}{Lemma}
\newcommand {\DNorm}[2]{ \mathchoice
    {\|\ee #1\ee\|\dd{#2}\,}
    {\| #1 \|_{#2}}
    {\| #1 \|_{#2}}
    {\| #1 \|_{#2}} }
\newcommand{\dd}[1]{_{\raise-0.6ex\hbox{$\scriptstyle #1$}}}
\newcommand{\ee}{\hskip0.15ex}
\newcommand{\sfE}{{\mathsf E}}
\newcommand{\sfe}{{\mathsf e}}
\newcommand{\sfu}{{\mathsf u}}
\title{Estimates for harmonic functions near pseudo-corners}
\author{Martin Costabel, Monique Dauge}
\date{Source: arXiv:2507.22178v2 (CC BY 4.0)}
\begin{document}
\maketitle

\noindent\emph{Source and licence.} Estimates for harmonic functions near pseudo-corners. Version arXiv:2507.22178v2 (CC BY 4.0), distributed by arXiv under the Creative Commons Attribution 4.0 International licence, http://creativecommons.org/licenses/by/4.0/, as recorded in the arXiv licence metadata for this exact version and shown on the abstract page https://arxiv.org/abs/2507.22178v2. Full licence notice: \texttt{licenses/arxiv-2507.22178-CC-BY-4.0.txt}.\par\smallskip

\noindent\textit{Editorial scope.} Counted unit: the article's lemma lem:S1 (source0.tex lines 969--977). The article's complete original proof of it is delivered with it (source0.tex lines 979--1021, one \texttt{proof} environment of 43 lines). No mathematics is written, repaired or re-derived: the statement and the proof are the article's own lines, in the article's own notation, and no retained line is substituted or re-flowed.  Retained uncounted as the source-local context the proof needs: lines 433--436; lines 714--719.  Nothing was omitted from any retained span, so the package carries no omission record.  No retained line contains a citation, so the wrapper carries no bibliography and no bibliography entry.  No retained line contains a figure environment, an image-inclusion command or drawing code, so no image is delivered and the package declares no asset dependency.  Editorial formatting only, in addition: an AMS \texttt{amsart} preamble that restates the article's own declarations and macros used by the retained lines, namely the environments \texttt{lemma} on separate counters, together with the standard packages those lines need.  The retained environments print in this document as Lemma 1 in order of appearance, whereas the article prints the same items with the numbering of its own full document.  The complete native source of the article as distributed in the arXiv e-print for this exact version is bundled unchanged as \texttt{sources/182396/source0.tex}.
\begin{equation}
\label{eq:lj}
  \lambda_j^\pm = 1-\frac{n}{2} \pm \sqrt{\left(1-\frac{n}{2}\right)^2 + \mu_j}\;,
\end{equation}

\begin{equation}
\label{eq:sfuepsb}
   \sum_{\sfe\in\sfE^\infty} \varepsilon^\sfe\,\DNorm{\sfu_\sfe}{H^1(\Omega)} 
   \le C \DNorm{f}{L^2(\Omega)}\quad
   \forall\varepsilon\in[0,\varepsilon_\star).
\end{equation}

\begin{lemma}
\label{lem:S1}
Let $\xi\in(0,\max\{1,\varepsilon_\star\})$. There exists $\gamma_1>0$ such that
\[
   \sum_{\sfe\in\sfE^\infty\!,\;\sfe\neq0} \varepsilon^{\sfe} \DNorm{\sfu_\sfe}{H^1(\Omega)}
   \le \gamma_1\, \varepsilon^{\lambda^+_1}
   {\DNorm{f}{L^2(\Omega)}},\quad\forall\varepsilon\in[0,\xi].
\]
\end{lemma}

\begin{proof}
Recall that $\sfE^\infty\setminus\{0\}$ is the set of all finite sums of the $\pm\lambda^\pm_j$ defined in \eqref{eq:lj}. The smallest element of $\{\pm\lambda^\pm_j,\;j\ge1\}$ is $\lambda^+_1$. Hence it is also the smallest element of $\sfE^\infty\setminus\{0\}$. Hence we can write
\begin{equation}
\label{eq:eneq0}
   \sum_{\sfe\in\sfE^\infty\!,\;\sfe\neq0} \varepsilon^{\sfe} \DNorm{\sfu_\sfe}{H^1(\Omega)} =
   \varepsilon^{\lambda^+_1} 
   \sum_{\sfe\in\sfE^\infty\setminus\{0\}} \varepsilon^{\sfe-\lambda^+_1} 
   \DNorm{\sfu_\sfe}{H^1(\Omega)}.
\end{equation}
Let $\delta>0$ be such that $\xi^{1-\delta}<\varepsilon_\star$. 
\begin{itemize}
\item There is a finite number of elements $\sfe\in\sfE^\infty\setminus\{0\}$ such that $\sfe<\delta^{-1}\lambda^+_1$.

\item For the other exponents $\sfe$, we have $\delta\sfe\ge \lambda^+_1$, hence if $\varepsilon\le1$
\[
   \varepsilon^{\sfe-\lambda^+_1} = 
   \varepsilon^{\sfe(1-\lambda^+_1/\sfe)} \le 
   \varepsilon^{\sfe(1-\delta)}=
   (\varepsilon^{1-\delta})^\sfe.
\]
\end{itemize}
Hence for $\varepsilon\le1$
\[
\begin{aligned}
   \sum_{\sfe\in\sfE^\infty\setminus\{0\}} \varepsilon^{\sfe-\lambda^+_1} 
   \DNorm{\sfu_\sfe}{H^1(\Omega)} &=
   \sum_{\lambda^+_1\le\sfe<\delta^{-1}\lambda^+_1}
   \varepsilon^{\sfe-\lambda^+_1} 
   \DNorm{\sfu_\sfe}{H^1(\Omega)} +
   \sum_{\sfe\ge\delta^{-1}\lambda^+_1}
   \varepsilon^{\sfe-\lambda^+_1} 
   \DNorm{\sfu_\sfe}{H^1(\Omega)} \\
&\le
   \sum_{\lambda^+_1\le\sfe<\delta^{-1}\lambda^+_1}
   \varepsilon^{\sfe-\lambda^+_1} 
   \DNorm{\sfu_\sfe}{H^1(\Omega)} +
   \sum_{\sfe\ge\delta^{-1}\lambda^+_1}
   (\varepsilon^{1-\delta})^\sfe 
   \DNorm{\sfu_\sfe}{H^1(\Omega)} .
\end{aligned}
\]
Therefore, if $\varepsilon\le\xi$, thanks to \eqref{eq:sfuepsb} the right-hand side defines a quantity bounded independently of $\varepsilon$, which, owing to \eqref{eq:eneq0}, proves the lemma.
\end{proof}

\end{document}
